\documentclass[a4paper,12pt]{article}
\usepackage[italian]{babel}
\usepackage[utf8]{inputenc}
\usepackage[T1]{fontenc}
\usepackage{geometry}
\geometry{margin=2.5cm}

\title{Verifica di Basi di Dati -- SQL}
\author{Prof. Fedeli Massimo -- Fabbrica Digitale}
\date{16/03/2026}

\begin{document}
	
	\maketitle
	
	\section*{Schema Logico}
	
	PAZIENTI(\underline{id\_paziente}, nome, cognome, data\_nascita, citt\`a)
	
	MEDICI(\underline{id\_medico}, nome, cognome, specializzazione)
	
	REPARTI(\underline{id\_reparto}, nome\_reparto, piano, id\_medico)
	
	RICOVERI(\underline{id\_paziente}, \underline{id\_reparto}, data\_ricovero)
	
	VISITE(\underline{id\_paziente}, \underline{id\_medico}, voto, data\_visita)
	
	
	\section*{Parte 1 -- Query SQL}
	
	Scrivere le query SQL per le seguenti richieste.
	
	\begin{enumerate}
		
\item Elencare nome, cognome e città di residenza dei pazienti
nati dopo il 1980, ordinati per data di nascita decrescente.

\item Elencare i reparti situati ai piani 2 o 3, mostrando
il nome del reparto, il piano e il numero totale di pazienti
attualmente ricoverati in ciascuno.

\item Elencare nome e cognome dei pazienti residenti a Roma
o Milano che hanno effettuato almeno una visita con
valutazione superiore a 7.
		
		\item Elencare i pazienti ordinati per cognome e nome.
		
		\item Elencare i reparti con il nome del medico responsabile.
		
		\item Elencare i pazienti che hanno effettuato almeno una visita.
		
		\item Elencare i pazienti che hanno ricevuto almeno una valutazione maggiore o uguale a 8.
		
		\item Calcolare la valutazione media per ogni reparto.
		
		\item Elencare i pazienti che hanno effettuato pi\`u di 2 visite.
		
		\item Elencare i medici che seguono pi\`u di un reparto.
		
		\item Elencare i reparti che non hanno ancora visite registrate.
		
		\item Trovare il paziente con la media delle valutazioni pi\`u alta.
		
		\item Elencare i pazienti ricoverati in tutti i reparti a cui sono assegnati.
		
		\item Elencare il numero di pazienti ricoverati in ciascun reparto.
		
		\item Trovare il reparto con la valutazione media pi\`u alta.
		
	\end{enumerate}
	
	
	\section*{Parte 2 -- INSERT}
	
	Inserire un nuovo paziente:
	Mario Rossi, nato il 15 marzo 2000, residente a Bologna.
	
	
	\section*{Parte 3 -- UPDATE}
	
	Aggiornare il piano del reparto ``Cardiologia'' portandolo al piano 3.
	

\section*{Parte 4 -- SQL DDL}

Scrivere le istruzioni SQL per creare le seguenti tabelle:

SALE(\underline{id\_sala}, edificio, capienza)

INTERVENTI(\underline{id\_intervento}, descrizione,
data\_intervento, id\_sala, id\_medico)

Vincoli:
\begin{itemize}
	\item \texttt{id\_sala} e \texttt{id\_intervento} chiavi primarie
	\item \texttt{capienza} $> 0$
	\item \texttt{id\_sala} in INTERVENTI chiave esterna verso SALE
	\item \texttt{id\_medico} in INTERVENTI chiave esterna verso MEDICI
	\item \texttt{data\_intervento} non può essere NULL
\end{itemize}
		
	\section*{Parte 5 -- Viste}
	
	Creare una vista chiamata \texttt{media\_pazienti} che mostri:
	\begin{itemize}
		\item \texttt{id\_paziente}
		\item \texttt{nome}
		\item \texttt{cognome}
		\item media delle valutazioni
	\end{itemize}
	
\end{document}